\documentclass[a4paper]{article}     
\usepackage{amsfonts}
\usepackage{amsmath}
\usepackage{amssymb}
\usepackage{graphicx}

\begin{document}

\noindent \large Auteur: Adsayan Gunaseelan \\
\noindent \large Studierichting: Informatica\\
\noindent \large Oefeningenreeks 1C nr. 31.5\\

\medskip

\normalsize

Opgave: $(\sqrt{6} + i\sqrt{2})^{12}$ Bereken in $\mathbb{C}$ door gebruik van goniometrische vormen\\

$\Rightarrow |z| = \sqrt{(\sqrt{6})^2+(\sqrt{2})^2} = \sqrt{6+2}=\sqrt{8}=2\sqrt{2}$\\

$\tan(\theta) = \dfrac{\sqrt{2}}{\sqrt{6}} = \dfrac{1}{\sqrt{3}}$\\

We weten dat $\tan\left(\dfrac{\pi}{6}\right) = \dfrac{1}{\sqrt{3}}$.\\

Dus $\theta = \dfrac{\pi}{6}$.\\


$\Rightarrow z^{12} = (2\sqrt{2})^{12}(\cos({12\cdot\dfrac{\pi}{6}})+i\sin({12\cdot\dfrac{\pi}{6}}))$\\

$=(2\sqrt{2})^{12}(\cos(2\pi)+i\sin(2\pi))$\\

$ =2^{\tfrac{3}{2}\cdot12}(\cos(2\pi)+i\sin(2\pi))$\\

$ =2^{18}(\cos(2\pi)+i\sin(2\pi))$\\

$ =2^{18}$\\

\begin{thebibliography}{999}
%\bibitem{a} Referentie
\end{thebibliography}
\end{document} 